\section{結論}\label{sec:conclusion}

本研究では、スライディングウィンドウに基づく密集区間マイニングに$(\varepsilon, \delta)$-差分プライバシー保証を付与するDP-DIMフレームワークを提案した。

主要な理論的貢献として、(1) ウィンドウカウントクエリの窓感度が$\Delta_W = 1$であることの形式的証明、(2) 有界隣接関係に基づくプライバシー定義の明確化、(3) ラプラスメカニズムおよびガウスメカニズムに基づくDP密集区間検出アルゴリズム、(4) 閾値安定性条件の導出、(5) SVTによる予算効率的な変種、(6) 多レベルApriori探索におけるプライバシー保証の正式な証明と高度合成定理による予算管理手法、(7) 滑らかな感度の分析による密集区間マイニングの構造的優位性の解明、(8) 情報理論的下界の導出によるラプラスメカニズムの準最適性の証明を示した。

実験結果より、以下が確認された：
\begin{itemize}
  \item $\varepsilon = 1.0$の実用的なプライバシーレベルで平均F1スコア0.60を達成し、非DP手法からの精度低下は6.5\%に留まること。$\varepsilon$に対する精度の平坦性（相対変動3.2\%）は、閾値からの距離に起因するDPノイズへの本質的頑健性として下界理論により説明される。
  \item ラプラスメカニズムがガウスメカニズムおよびSVTと比較して最も高い精度を示すこと。非DPオラクルとの定量的比較により精度低下の程度を明確にした。
  \item 実行時間がトランザクション数に対してほぼ線形にスケールし、DPオーバーヘッドが無視できること。
  \item 高度合成定理により$k = 500$クエリで逐次合成比77.5\%の予算削減が可能であること。
  \item パラメータ感度分析により、$\theta$と$W$の増大がDP精度を改善するメカニズムを解明し、$\theta \geq 2\gamma$の設計指針を導出したこと。
\end{itemize}

\subsection{今後の課題}

\textbf{ユーザーレベル差分プライバシー}：本研究ではトランザクションレベルのプライバシー（1トランザクションの置換を保護）を採用したが、実応用では1ユーザーが複数のトランザクションを生成する。ユーザーレベルDPでは、1ユーザーの全トランザクション（最大$B$個）の置換を保護する必要があり、窓感度は$\Delta_W = \min(B, W+1)$に増大する。これはグループプライバシー~\cite{dwork2014algorithmic}の枠組みで$\varepsilon$を$B$倍にすることで対応可能であるが、精度の大幅な低下が予想されるため、ユーザーごとのサンプリングやサブサンプリング増幅~\cite{balle2018privacy}との組み合わせが有望である。

\textbf{実データでの評価}：合成データに加え、小売（Dunnhumby, Online Retail）・医療・ウェブログなどの実データでの評価が必要である。特に、密集パターンの疎密度やアイテム数がDP精度に与える影響を調査する。実データではノイズ特性が合成データと異なる可能性があり、外部妥当性の確認が不可欠である。

\textbf{R\'{e}nyi差分プライバシー}：RDP~\cite{mironov2017renyi}およびゼロ集中差分プライバシー（zCDP）~\cite{bun2016concentrated}を用いた合成は高度合成よりも厳密な上界を与える場合があり、特にガウスメカニズムとの組み合わせで予算管理のさらなる改善が期待される。

\textbf{適応的予算配分}：候補アイテムセットの数やウィンドウ位置に応じた適応的な予算配分戦略の開発により、固定配分を上回る精度が達成できる可能性がある。Aprioriの枝刈り情報を活用し、有望な候補により多くの予算を配分する戦略が考えられる。

\textbf{局所差分プライバシー}：本研究は信頼できるデータ管理者が存在する中央集権型DPを想定している。各ユーザーが自身のトランザクションにローカルにノイズを加える局所差分プライバシー（LDP）モデルへの拡張も重要な方向性である。LDPではサーバーへの信頼が不要であるが、精度の大幅な低下が見込まれるため、ランダム化応答やUE（Unary Encoding）プロトコルとの統合が課題となる。
